
\documentclass{article}
\usepackage{colortbl}
\usepackage{makecell}
\usepackage{multirow}
\usepackage{supertabular}

\begin{document}

\newcounter{utterance}

\centering \large Interaction Transcript for game `dond', experiment `coop\_it', episode 3 with gpt{-}5{-}mini{-}2025{-}08{-}07{-}no{-}reasoning{-}t1.0.
\vspace{24pt}

{ \footnotesize  \setcounter{utterance}{1}
\setlength{\tabcolsep}{0pt}
\begin{supertabular}{c@{$\;$}|p{.15\linewidth}@{}p{.15\linewidth}p{.15\linewidth}p{.15\linewidth}p{.15\linewidth}p{.15\linewidth}}
    \# & $\;$A & \multicolumn{4}{c}{Game Master} & $\;\:$B\\
    \hline

    \theutterance \stepcounter{utterance}  
    & & \multicolumn{4}{p{0.6\linewidth}}{
        \cellcolor[rgb]{0.9,0.9,0.9}{
            \makecell[{{p{\linewidth}}}]{
                \texttt{\tiny{[P1$\langle$GM]}}
                \texttt{State giocando a un gioco di negoziazione cooperativa in cui dovete accordarvi su come dividere una serie di oggetti tra voi e un altro giocatore.} \\
\\ 
\texttt{Regole:} \\
\texttt{(a) A Lei e all'altro giocatore viene dato un insieme di oggetti. Ognuno di voi riceve anche una funzione di valore segreta, che rappresenta il valore di ciascun tipo di oggetto.} \\
\texttt{(b) Si scambiano messaggi con l'altro giocatore per concordare chi si aggiudica gli oggetti. Potete inviare un massimo di 5 messaggi ciascuno, oppure terminare in anticipo facendo una proposta segreta in qualsiasi momento.} \\
\texttt{(c) A ciascuno di voi viene chiesto di inviare una proposta segreta indicando gli oggetti che desiderate, formattata tra parentesi quadre come segue: "{[}Proposta: <numero> <nome oggetto>, <numero> <nome oggetto>, <...>{]}".} \\
\texttt{(d) Se le vostre proposte sono complementari, cioè ci sono abbastanza oggetti per soddisfare entrambe le proposte, a ciascun giocatore viene assegnato un punteggio basato sulla somma dei valori degli oggetti ricevuti. In caso contrario, entrambi ricevono zero punti.} \\
\texttt{(e) L'obiettivo è massimizzare la somma del proprio punteggio e di quello dell'altro giocatore.} \\
\\ 
\texttt{Cominciamo.} \\
\\ 
\texttt{L'insieme degli oggetti disponibili è:} \\
\\ 
\texttt{1 carota, 4 borse, 1 cipolla} \\
\\ 
\texttt{La funzione valore segreta è:} \\
\\ 
\texttt{carota: 3, borsa: 1, cipolla: 3} \\
\\ 
\texttt{IMPORTANTE: i vostri messaggi, a meno che non si tratti di una proposta segreta, vengono trasmessi direttamente all'altro giocatore, quindi non includete alcuna risposta alle regole o testo di annuncio del vostro messaggio. Per fare una proposta segreta, utilizzate il formato indicato. Non utilizzare delle parentesi quadre quando si comunica all'altro giocatore, altrimenti verrà interpretata come una proposta segreta.} \\
            }
        }
    }
    & \\ \\

    \theutterance \stepcounter{utterance}  
    & \multicolumn{4}{p{0.6\linewidth}}{
        \cellcolor[rgb]{0.9,0.9,0.9}{
            \makecell[{{p{\linewidth}}}]{
                \texttt{\tiny{[P1$\rangle$GM]}}
                \texttt{Ciao! Condivido volentieri le mie preferenze per facilitare un accordo efficiente: per me carota vale 3, borsa 1 ciascuna, cipolla 3. Quindi carota e cipolla sono molto importanti, le borse hanno valore minore ma comunque positivo.} \\
\\ 
\texttt{Proposta iniziale non vincolante: io prendo carota e cipolla, tu prendi le 4 borse. Questo mi darebbe 6 e a te 4. Se per te una tra carota o cipolla vale molto, possiamo scambiare: ad esempio tu prendi cipolla e 2 borse, io prendo carota e 2 borse. Fammi sapere i tuoi valori o preferenze così ottimizziamo la somma totale.} \\
            }
        }
    }
    & & \\ \\

    \theutterance \stepcounter{utterance}  
    & & & \multicolumn{2}{p{0.3\linewidth}}{
        \cellcolor[rgb]{0.9,0.9,0.9}{
            \makecell[{{p{\linewidth}}}]{
                \texttt{\tiny{[GM$|$GM]}}
                \texttt{continue} \\
            }
        }
    }
    & & \\ \\

    \theutterance \stepcounter{utterance}  
    & & & \multicolumn{2}{p{0.3\linewidth}}{
        \cellcolor[rgb]{0.9,0.9,0.9}{
            \makecell[{{p{\linewidth}}}]{
                \texttt{\tiny{[GM$|$GM]}}
                \texttt{Ciao! Condivido volentieri le mie preferenze per facilitare un accordo efficiente: per me carota vale 3, borsa 1 ciascuna, cipolla 3. Quindi carota e cipolla sono molto importanti, le borse hanno valore minore ma comunque positivo.} \\
\\ 
\texttt{Proposta iniziale non vincolante: io prendo carota e cipolla, tu prendi le 4 borse. Questo mi darebbe 6 e a te 4. Se per te una tra carota o cipolla vale molto, possiamo scambiare: ad esempio tu prendi cipolla e 2 borse, io prendo carota e 2 borse. Fammi sapere i tuoi valori o preferenze così ottimizziamo la somma totale.} \\
            }
        }
    }
    & & \\ \\

    \theutterance \stepcounter{utterance}  
    & & \multicolumn{4}{p{0.6\linewidth}}{
        \cellcolor[rgb]{0.9,0.9,0.9}{
            \makecell[{{p{\linewidth}}}]{
                \texttt{\tiny{[GM$\rangle$P2]}}
                \texttt{State giocando a un gioco di negoziazione cooperativa in cui dovete accordarvi su come dividere una serie di oggetti tra voi e un altro giocatore.} \\
\\ 
\texttt{Regole:} \\
\texttt{(a) A Lei e all'altro giocatore viene dato un insieme di oggetti. Ognuno di voi riceve anche una funzione di valore segreta, che rappresenta il valore di ciascun tipo di oggetto.} \\
\texttt{(b) Si scambiano messaggi con l'altro giocatore per concordare chi si aggiudica gli oggetti. Potete inviare un massimo di 5 messaggi ciascuno, oppure terminare in anticipo facendo una proposta segreta in qualsiasi momento.} \\
\texttt{(c) A ciascuno di voi viene chiesto di inviare una proposta segreta indicando gli oggetti che desiderate, formattata tra parentesi quadre come segue: "{[}Proposta: <numero> <nome oggetto>, <numero> <nome oggetto>, <...>{]}".} \\
\texttt{(d) Se le vostre proposte sono complementari, cioè ci sono abbastanza oggetti per soddisfare entrambe le proposte, a ciascun giocatore viene assegnato un punteggio basato sulla somma dei valori degli oggetti ricevuti. In caso contrario, entrambi ricevono zero punti.} \\
\texttt{(e) L'obiettivo è massimizzare la somma del proprio punteggio e di quello dell'altro giocatore.} \\
\\ 
\texttt{Cominciamo.} \\
\\ 
\texttt{L'insieme degli oggetti disponibili è:} \\
\\ 
\texttt{1 carota, 4 borse, 1 cipolla} \\
\\ 
\texttt{La funzione valore segreta è:} \\
\\ 
\texttt{carota: 0, borsa: 1, cipolla: 6} \\
\\ 
\texttt{IMPORTANTE: i vostri messaggi, a meno che non si tratti di una proposta segreta, vengono trasmessi direttamente all'altro giocatore, quindi non includete alcuna risposta alle regole o testo di annuncio del vostro messaggio. Per fare una proposta segreta, utilizzate il formato indicato. Non utilizzare delle parentesi quadre quando si comunica all'altro giocatore, altrimenti verrà interpretata come una proposta segreta.} \\
\\ 
\\ 
\texttt{Ciao! Condivido volentieri le mie preferenze per facilitare un accordo efficiente: per me carota vale 3, borsa 1 ciascuna, cipolla 3. Quindi carota e cipolla sono molto importanti, le borse hanno valore minore ma comunque positivo.} \\
\\ 
\texttt{Proposta iniziale non vincolante: io prendo carota e cipolla, tu prendi le 4 borse. Questo mi darebbe 6 e a te 4. Se per te una tra carota o cipolla vale molto, possiamo scambiare: ad esempio tu prendi cipolla e 2 borse, io prendo carota e 2 borse. Fammi sapere i tuoi valori o preferenze così ottimizziamo la somma totale.} \\
            }
        }
    }
    & \\ \\

    \theutterance \stepcounter{utterance}  
    & & & \multicolumn{4}{p{0.6\linewidth}}{
        \cellcolor[rgb]{0.9,0.9,0.9}{
            \makecell[{{p{\linewidth}}}]{
                \texttt{\tiny{[GM$\langle$P2]}}
                \texttt{Grazie per la trasparenza. Con i miei valori, la cipolla vale tantissimo (6), le borse valgono 1 ciascuna e la carota 0. Quindi massimizzare la somma totale richiede che la cipolla vada a me, mentre le borse possono andare a chi ne trae più valore marginale a parità.} \\
\\ 
\texttt{Dato che per te carota=3 e cipolla=3, una divisione che massimizza la somma è:} \\
\texttt{{-} Io prendo la cipolla (valore per me 6)} \\
\texttt{{-} Tu prendi la carota (3 per te) e le 4 borse (4 per te)} \\
\\ 
\texttt{Totale somma = 6 + (3+4) = 13, che è superiore alla tua proposta iniziale (10).} \\
\\ 
\texttt{Se sei d’accordo, confermiamo. Se preferisci una variante, posso cedere 1 borsa a me o a te non cambia il totale per me, ma per te sì: tu ne trai 1 punto per borsa, io pure 1 per borsa. Quindi tenere tutte le borse a te massimizza anche il tuo punteggio senza ridurre la somma.} \\
            }
        }
    }
     \\ \\

    \theutterance \stepcounter{utterance}  
    & & & \multicolumn{2}{p{0.3\linewidth}}{
        \cellcolor[rgb]{0.9,0.9,0.9}{
            \makecell[{{p{\linewidth}}}]{
                \texttt{\tiny{[GM$|$GM]}}
                \texttt{continue} \\
            }
        }
    }
    & & \\ \\

    \theutterance \stepcounter{utterance}  
    & & & \multicolumn{2}{p{0.3\linewidth}}{
        \cellcolor[rgb]{0.9,0.9,0.9}{
            \makecell[{{p{\linewidth}}}]{
                \texttt{\tiny{[GM$|$GM]}}
                \texttt{Grazie per la trasparenza. Con i miei valori, la cipolla vale tantissimo (6), le borse valgono 1 ciascuna e la carota 0. Quindi massimizzare la somma totale richiede che la cipolla vada a me, mentre le borse possono andare a chi ne trae più valore marginale a parità.} \\
\\ 
\texttt{Dato che per te carota=3 e cipolla=3, una divisione che massimizza la somma è:} \\
\texttt{{-} Io prendo la cipolla (valore per me 6)} \\
\texttt{{-} Tu prendi la carota (3 per te) e le 4 borse (4 per te)} \\
\\ 
\texttt{Totale somma = 6 + (3+4) = 13, che è superiore alla tua proposta iniziale (10).} \\
\\ 
\texttt{Se sei d’accordo, confermiamo. Se preferisci una variante, posso cedere 1 borsa a me o a te non cambia il totale per me, ma per te sì: tu ne trai 1 punto per borsa, io pure 1 per borsa. Quindi tenere tutte le borse a te massimizza anche il tuo punteggio senza ridurre la somma.} \\
            }
        }
    }
    & & \\ \\

    \theutterance \stepcounter{utterance}  
    & & \multicolumn{4}{p{0.6\linewidth}}{
        \cellcolor[rgb]{0.9,0.9,0.9}{
            \makecell[{{p{\linewidth}}}]{
                \texttt{\tiny{[P1$\langle$GM]}}
                \texttt{Grazie per la trasparenza. Con i miei valori, la cipolla vale tantissimo (6), le borse valgono 1 ciascuna e la carota 0. Quindi massimizzare la somma totale richiede che la cipolla vada a me, mentre le borse possono andare a chi ne trae più valore marginale a parità.} \\
\\ 
\texttt{Dato che per te carota=3 e cipolla=3, una divisione che massimizza la somma è:} \\
\texttt{{-} Io prendo la cipolla (valore per me 6)} \\
\texttt{{-} Tu prendi la carota (3 per te) e le 4 borse (4 per te)} \\
\\ 
\texttt{Totale somma = 6 + (3+4) = 13, che è superiore alla tua proposta iniziale (10).} \\
\\ 
\texttt{Se sei d’accordo, confermiamo. Se preferisci una variante, posso cedere 1 borsa a me o a te non cambia il totale per me, ma per te sì: tu ne trai 1 punto per borsa, io pure 1 per borsa. Quindi tenere tutte le borse a te massimizza anche il tuo punteggio senza ridurre la somma.} \\
            }
        }
    }
    & \\ \\

    \theutterance \stepcounter{utterance}  
    & \multicolumn{4}{p{0.6\linewidth}}{
        \cellcolor[rgb]{0.9,0.9,0.9}{
            \makecell[{{p{\linewidth}}}]{
                \texttt{\tiny{[P1$\rangle$GM]}}
                \texttt{Concordo: cipolla a te; a me carota e 4 borse. Procedo con la proposta segreta.} \\
\\ 
\texttt{{[}Proposta: 1 carota, 4 borse{]}} \\
            }
        }
    }
    & & \\ \\

    \theutterance \stepcounter{utterance}  
    & & & \multicolumn{2}{p{0.3\linewidth}}{
        \cellcolor[rgb]{0.9,0.9,0.9}{
            \makecell[{{p{\linewidth}}}]{
                \texttt{\tiny{[GM$|$GM]}}
                \texttt{proposal} \\
            }
        }
    }
    & & \\ \\

    \theutterance \stepcounter{utterance}  
    & & & \multicolumn{2}{p{0.3\linewidth}}{
        \cellcolor[rgb]{0.9,0.9,0.9}{
            \makecell[{{p{\linewidth}}}]{
                \texttt{\tiny{[GM$|$GM]}}
                [1, 4, 0]
            }
        }
    }
    & & \\ \\

    \theutterance \stepcounter{utterance}  
    & & \multicolumn{4}{p{0.6\linewidth}}{
        \cellcolor[rgb]{0.9,0.9,0.9}{
            \makecell[{{p{\linewidth}}}]{
                \texttt{\tiny{[GM$\rangle$P2]}}
                \texttt{L'altro giocatore ha presentato una proposta segreta. Ora tocca a Lei presentare una sua proposta.} \\
            }
        }
    }
    & \\ \\

    \theutterance \stepcounter{utterance}  
    & & & \multicolumn{4}{p{0.6\linewidth}}{
        \cellcolor[rgb]{0.9,0.9,0.9}{
            \makecell[{{p{\linewidth}}}]{
                \texttt{\tiny{[GM$\langle$P2]}}
                \texttt{{[}Proposta: 1 cipolla, 4 borse, 0 carota{]}} \\
            }
        }
    }
     \\ \\

    \theutterance \stepcounter{utterance}  
    & & & \multicolumn{2}{p{0.3\linewidth}}{
        \cellcolor[rgb]{0.9,0.9,0.9}{
            \makecell[{{p{\linewidth}}}]{
                \texttt{\tiny{[GM$|$GM]}}
                \texttt{proposal} \\
            }
        }
    }
    & & \\ \\

    \theutterance \stepcounter{utterance}  
    & & & \multicolumn{2}{p{0.3\linewidth}}{
        \cellcolor[rgb]{0.9,0.9,0.9}{
            \makecell[{{p{\linewidth}}}]{
                \texttt{\tiny{[GM$|$GM]}}
                [0, 4, 1]
            }
        }
    }
    & & \\ \\

    \theutterance \stepcounter{utterance}  
    & & & \multicolumn{2}{p{0.3\linewidth}}{
        \cellcolor[rgb]{0.9,0.9,0.9}{
            \makecell[{{p{\linewidth}}}]{
                \texttt{\tiny{[GM$|$GM]}}
                [[1, 4, 0], [0, 4, 1]]
            }
        }
    }
    & & \\ \\

\end{supertabular}
}

\end{document}
